\paragraph{Auxiliary internal-routing objectives.}
The rollout coefficient loss specifies the prediction term. Learned
internal group-routing configurations additionally use group-balance,
entropy, and Reynolds-partition supervision terms. These concern
expert-group selection within a specialist, not the outer E2 router
or the T2-C gate. Their exact signs, reductions, and per-run weights
are not established by the architecture report; no universal
auxiliary-loss formula or coefficient set is therefore asserted here.

\begin{table}[t]
\centering
\small
\caption{Scope of internal group routing. Deployment policy alone does
not establish the auxiliary objectives used during training.}
\label{tab:internal_objective_scope}
\begin{tabular}{@{}lp{0.30\linewidth}p{0.40\linewidth}@{}}
\toprule
Instance & Group configuration & Interpretation\\
\midrule
Circular $H$ & Single group; group router frozen & No learned group choice; channel Top-2 remains distinct.\\
Circular $P$ & Three groups; learned Top-1 & Straight-through group selection with auxiliary group objectives.\\
Square $S$ & Three groups; learned Top-1 & Group routing is distinct from frozen velocity expert weights.\\
Square $H$ & Single group; group router frozen & No learned group choice; data-only overlap output path.\\
Square $P$ & Three stored groups; fixed deployment index 1 & Training auxiliary weights cannot be inferred from fixed deployment.\\
Pinball $S/H/P$ & Deep-FNN-H3; no group router & Internal group-routing objectives are not applicable.\\
\bottomrule
\end{tabular}
\end{table}

For learned group selection, the straight-through gate has a hard
one-hot forward value and a softmax backward path. Source confirmation
of this computation is not a separate dynamic measurement of nonzero
router gradients. Frozen expert parameters also do not imply constant
corrections: their inputs and mixture weights can still vary.
